\documentclass[10pt,a4paper]{amsart}
\usepackage[margin=19mm]{geometry}
\usepackage{amsmath,amssymb,mathtools}
\usepackage[scr=rsfs,cal=euler]{mathalfa}
\usepackage{xfrac}
\usepackage[unicode,hidelinks,bookmarks=false]{hyperref}
\newcommand{\ie}{i.e.\ }
\newcommand{\cf}{cf.\ }
\newcommand{\resp}{respectively\ }
\newcommand{\Subsection}{\subsection}
\providecommand{\nobreakdash}{\nobreak\hbox{-}}
\setlength{\emergencystretch}{2em}
\multlinegap0pt
\usepackage{bm}
\usepackage{bbm}
\usepackage{dsfont}
\usepackage{xspace}
\usepackage{cancel}
\usepackage{mathtools}
\usepackage[shortlabels]{enumitem}
\usepackage{amsthm}
\makeatletter
\@ifundefined{equation}{\newtheorem{equation}{Equation}}{}
\@ifundefined{lemm}{\newtheorem{lemm}[equation]{Lemma}}{}
\makeatother
\newcommand{\M}{\mathrm{M}}
\newcommand{\Q}{\mathbf{Q}}
\newcommand{\Tr}{\mathrm{Tr}}
\newcommand{\Z}{\mathbf{Z}}
\title[On families of strongly divisible modules of rank 2]{On families of strongly divisible modules of rank 2}
\author{Seongjae Han, Chol Park}
\date{Source: arXiv:2503.03994v2 (CC BY 4.0)}
\begin{document}
\maketitle

\noindent\emph{Source and licence.} On families of strongly divisible modules of rank 2. Version arXiv:2503.03994v2 (CC BY 4.0), distributed under the Creative Commons Attribution 4.0 International licence, as recorded in the arXiv licence metadata for this exact version and shown on the abstract page https://arxiv.org/abs/2503.03994v2. Full rights notice: licenses/arxiv-2503.03994-CC-BY-4.0.txt.\par\smallskip

\noindent\textit{Editorial scope.} Counted unit: the Lemm whose source label is \texttt{lemma: m+1 to r} in the article (source0.tex lines 1662--1682, source label \texttt{lemma: m+1 to r}), delivered together with the article's own complete original proof of it (source0.tex lines 1804--2221, one \texttt{proof} environment of 418 lines). No mathematics is written, repaired or re-derived: statement and proof are the article's own text line for line. The proof is detached from the statement in the source: the article states the result at lines 1662--1682 and prints its proof at lines 1804--2221, and the proof environment's own header names the statement, so the association is the article's own. Required context retained uncounted: lemma: k=0,k=r (lines 1200--1211); lemma: PQdegree (lines 1253--1260); eq: inverse of B\_k,x (lines 1704--1706); lemma-QAinv (lines 1721--1728). Every definition, display and label of those lines is retained unchanged. Omitted from this package and not redistributed: 1--1199, 1212--1252, 1261--1661, 1683--1703, 1707--1720, 1729--1803, 2222--9214. Nothing inside a retained excerpt is omitted or altered: every retained body is delivered exactly as the article's lines are written, so no omission receipt of any kind accompanies this package. Citations: the retained excerpts carry no citation command at all, so no bibliography entry of the article is transcribed and no reference is redistributed. Reference adaptation: none was needed and none was made; every reference inside the delivered unit resolves inside it, so no unresolved reference or citation is produced. Printed slips kept literal: no printed word, symbol, hypothesis or number of the counted statement or of its proof was altered. Layout and class: the article's own class is \texttt{amsart} with packages including enumitem, amsmath, latexsym, amsfonts, amssymb, graphicx, mathrsfs, graphpap, color, geometry, multirow, hyperref, xy; the wrapper is the lane's pinned \texttt{amsart} editorial wrapper at 10pt on A4, which restates the theorem-like environments the retained text needs and adds the article's own macro definitions that the retained text needs (\texttt{\textbackslash M}, \texttt{\textbackslash Q}, \texttt{\textbackslash Tr}, \texttt{\textbackslash Z}), with extra wrapper packages (bm, bbm, dsfont, xspace, cancel, mathtools, enumitem). The delivered numbering is the unit's own. Assets: no retained span contains a figure environment, a graphics-inclusion command, a PDF-page inclusion, a table or a TikZ picture; no image file is copied into this package, the package declares no asset dependency and no image file is redistributed with it. Licence: version arXiv:2503.03994v2 of this article is distributed under the Creative Commons Attribution 4.0 International licence, as recorded in the arXiv licence metadata for this exact version and shown on the abstract page https://arxiv.org/abs/2503.03994v2. A full rights notice is delivered at licenses/arxiv-2503.03994-CC-BY-4.0.txt. Characters: every retained excerpt is pure ASCII, so the delivered engine, XeLaTeX with its default Latin Modern text font, reports no missing character.
\begin{lemm}\label{lemma: k=0,k=r}
For $0\leq i\leq r'-1$ we have
\begin{equation*}
   Q_{0,i}(T,x)=[T^i(x+Tf_{r'}(T))]_{T}^{(r'-1)},\quad  P_{0,i}(T,x)=T^i=Q_{r',i}(T,x),
\end{equation*}
and
\begin{equation*}
    P_{r',i}(T,x) =  x^{-r'}\sum_{s=i}^{r'-1}x^s (-1)^{r'-s-1}\Big[T^{i+r'-s-1}f_{r'}(T)^{r'-s-1}\Big]_{T}^{(r'-1)}.
\end{equation*}
In particular, $P_{k,i}(T,x)$ and $Q_{k,i}(T,x)$ belong to $\Z_{(p)}[x,\frac{1}{x},T]$ if $k=0$ or $k=r'$, and
$$Q_{k,i}(T,x)=[(x+Tf_{r'}(T))P_{k,i}(T,x)]_{T}^{(r'-1)}$$ for $k\in\{0,r'\}$ and for all $0\leq i\leq r'-1$.
\end{lemm}

\begin{lemm}\label{lemma: PQdegree}
Let $1\leq k\leq r'-1$, and assume that $x\in E$ with $\det (B_{k,x})\neq 0$. Then for each $0\leq i\leq r'-1$, $$\big(P_{k,i}(T,x),Q_{k,i}(T,x)\big)$$ is the unique pair of polynomials in $T$ satisfying the following properties:
    \begin{enumerate}[leftmargin=*]
        \item $Q_{k,i}(T,x)=[(x+Tf_{r'}(T))P_{k,i}(T,x)]_{T}^{(r'-1)}$;
        \item if $0\leq i\leq k-1$, then $\deg_T P_{k,i}(T,x)\leq k-1$ and $\deg_T\big( Q_{k,i}(T,x)-T^{i+r'-k} \big)\leq r'-k-1$;
        \item if $k\leq i\leq r'-1$, then $\deg_T\big( P_{k,i}(T,x)-T^i \big)\leq k-1$ and $\deg_T Q_{k,i}(T,x)\leq r'-k-1$.
    \end{enumerate}
\end{lemm}

\begin{equation}\label{eq: inverse of B_k,x}
    B_{k,x}^{-1}=B_{k,0}^{-1}-xB_{k,0}^{-1}U_t(I_t+xV_tB_{k,0}^{-1}U_t)^{-1}V_tB_{k,0}^{-1},
\end{equation}

\begin{lemm}\label{lemma-QAinv}
Let $m+1\leq k<r'$. Then
    \begin{enumerate}[leftmargin=*]
      \item $A'_k$ is invertible and $(A'_k)^{-1}\in\M_{t\times t}(\Z_{(p)})$;
      \item $B_{r'-k,0}^{-1}=A''_k-A'''_k(A'_k)^{-1}A_k$;
      \item $d_k(x)=\det(xI_t+(A'_k)^{-1})\in \Z_{(p)}[x]$, which is monic of degree $t$.
    \end{enumerate}
\end{lemm}

\begin{lemm}\label{lemma: m+1 to r}
    Let $m+1\leq k\leq r'$. Then
    \begin{equation*}
       t_{k,i}-1= s_{k,i}=
        \begin{cases}
            2k-r'-1 & \mbox{if } 0\leq i\leq 2k-r'-1; \\
            2k-r' & \mbox{if } 2k-r'\leq i\leq k-1;\\
            2k-r'+1 & \mbox{if } k\leq i\leq r'-1,
        \end{cases}
    \end{equation*}
    and
    \begin{equation*}
       Q_{k,i,t_{k,i}}(T)= P_{k,i,s_{k,i}}(T)
        =
        \begin{cases}
            P_{r'-k,i+r'-k}(T,0) & \mbox{if } 0\leq i\leq 2k-r'-1; \\
            P_{r'-k, i+r'-2k}(T,0) & \mbox{if } 2k-r'\leq i\leq k-1;\\
            -P_{r'-k,i-k}(T,0) & \mbox{if } k\leq i\leq r'-1.
        \end{cases}
    \end{equation*}
\end{lemm}

\begin{proof}[Proof of Lemma~\ref{lemma: m+1 to r}]
Let $t=2k-r'$, and we first consider the case $k=r'$. But this case immediately follows from Lemma~\ref{lemma: k=0,k=r}.




From now on, we assume that $m+1\leq k\leq r'-1$. We have the following identities:
$$B_{k,0}^{-1}U_t
=
        \begin{bmatrix}
          A'_k \\
          A'''_k
        \end{bmatrix},
\qquad
    V_tB_{k,0}^{-1}
        =
        \begin{bmatrix}
          A_k & A'_k
        \end{bmatrix},
\qquad\mbox{and}\qquad
    V_tB_{k,0}^{-1}U_t
        =A'_k.$$
We consider $y=x^{-1}$ to compute the terms with highest power of $x$. As $A'_k$ is invertible by Lemma~\ref{lemma-QAinv}~(i), we observe that
    \begin{align*}
        (I_t+xV_tB_{k,0}^{-1}U_t)^{-1}
        =
        (I_t+xA'_k)^{-1}
        =y(A'_k)^{-1}(y(A'_k)^{-1}+I_t)^{-1}.
    \end{align*}
    On $\Q[[y]]$, we have
    \begin{align*}
        (y(A'_k)^{-1}+I_t)^{-1}
        =\sum_{n=0}^\infty (y(A'_k)^{-1})^n
        \in \M_{t\times t}\big( \Q[[y]] \big),
    \end{align*}
    which gives
    \begin{align*}
        (I_t+xA'_k)^{-1}
        \equiv
        y(A'_k)^{-1}+y^2((A'_k)^{-1})^2 \pmod{y^3}.
    \end{align*}
Combining this to \eqref{eq: inverse of B_k,x} gives
    \begin{align*}
        B_{k,x}^{-1}
        &=
        B_{k,0}^{-1}
        -x
        \begin{bmatrix}
            A'_k \\
            A'''_k
        \end{bmatrix}
        (I_t+xA'_k)^{-1}
        \begin{bmatrix}
            A_k & A'_k
        \end{bmatrix}
        \\
        &\equiv
        B_{k,0}^{-1}
        -
        \begin{bmatrix}
            A'_k \\
            A'''_k
        \end{bmatrix}
        \big((A'_k)^{-1}+y((A'_k)^{-1})^2\big)
        \begin{bmatrix}
            A_k & A'_k
        \end{bmatrix}
        \pmod{y^2}.
    \end{align*}

We set
    \begin{align*}
        \tilde A
        :=
        \begin{bmatrix}
            A'_k \\
            A'''_k
        \end{bmatrix}
        (A'_k)^{-1}
        \begin{bmatrix}
            A_k & A'_k
        \end{bmatrix}
        \qquad\mbox{and}\qquad
        \tilde A'
        :=
        \begin{bmatrix}
            A'_k \\
            A'''_k
        \end{bmatrix}
        ((A'_k)^{-1})^2
        \begin{bmatrix}
            A_k & A'_k
        \end{bmatrix}
    \end{align*}
so that we have $B_{k,x}^{-1}\equiv B_{k,0}^{-1}-\tilde A-y\tilde A' \pmod{y^2}$. Then we have
    \begin{align*}
        \tilde A
        =
        \begin{bmatrix}
            A_k & A'_k \\
            A'''_k(A'_k)^{-1}A_k & A'''_k
        \end{bmatrix}.
    \end{align*}
By Lemma~\ref{lemma-QAinv}~(ii),
    \begin{equation}\label{eq: inverse of B_k,x-tilde A}
        B_{k,0}^{-1}-\tilde A
        =
        \begin{bmatrix}
            0_{t,k-t} & 0_{t,t} \\
            A''_k-A'''_k(A'_k)^{-1}A_k & 0_{k-t,t}
        \end{bmatrix}
        =
        \begin{bmatrix}
            0_{t,k-t} & 0_{t,t} \\
            B_{k-t,0}^{-1} & 0_{k-t,t}
        \end{bmatrix}.
    \end{equation}
For $\tilde A'$, we observe that
    \begin{equation}\label{eq: tilde A prime}
        \tilde A'
        =B_{k,0}^{-1}U_t((A'_k)^{-1})^2V_tB_{k,0}^{-1}=
        \begin{bmatrix}
            (A'_k)^{-1}A_k & I_t \\
            A'''_k((A'_k)^{-1})^2A_k & A'''_k(A'_k)^{-1}
        \end{bmatrix}.
    \end{equation}
Then it is easy to see that
    \begin{align*}
        M_{r',0}
        \begin{bmatrix}
            0_{r'-k,t} \\
            \tilde A'
        \end{bmatrix}
        =
        \begin{bmatrix}
            A_{k,0} & B_{k,0} \\
            C_{k,0} & D_{k,0}
        \end{bmatrix}
        \begin{bmatrix}
            0_{r'-k,t} \\
            \tilde A'
        \end{bmatrix}
        =
        \begin{bmatrix}
          B_{k,0}\tilde A' \\
          D_{k,0}\tilde A'
        \end{bmatrix}.
    \end{align*}
On the upper block, from the first equality of \eqref{eq: tilde A prime} we observe
    \begin{align*}
        B_{k,0}\tilde{A}'
        =
        \begin{bmatrix}
            0_{k-t,t} \\
            I_t
        \end{bmatrix}
        ((A'_k)^{-1})^2V_tB_{k,0}^{-1}
        =
        \begin{bmatrix}
            0_{k-t,k} \\
            ((A'_k)^{-1})^2V_tB_{k,0}^{-1}
        \end{bmatrix}.
    \end{align*}




Now we compute $P_{k,i,s_{k,i}}(T)$ for $0\leq i\leq k-1$. Recall that $\mathbf{e}'_1,\cdots,\mathbf{e}'_k$ be the standard basis of $\Q^k$, and that we have
    \begin{align*}
        P_{k,i}(T,x)
        =
        \begin{bmatrix}
            T^{k-1} & \cdots & 1
        \end{bmatrix}
        B_{k,x}^{-1}\mathbf{e}'_{k-i}.
    \end{align*}
From Lemma~\ref{lemma-QAinv}~(iii), we have
    \begin{align*}
        d_k(x)
       =\det(y^{-1}I_t+(A'_k)^{-1})
       \equiv
        y^{-t}+y^{-t+1}\Tr((A'_k)^{-1})
        \pmod{y^{-t+2}},
    \end{align*}
and so
    \begin{align*}
        \tilde P_{k,i}(T,x)
        &=
        \sum_{s=0}^{s_{k,i}}y^{-s} P_{k,i,s}(T)
        =d_k(x)P_{k,i}(T,x)
        \\&\equiv
        \big(y^{-t}+y^{-t+1}\Tr((A'_k)^{-1})\big)
        \begin{bmatrix}
            T^{k-1} & \cdots & 1
        \end{bmatrix}
        (B_{k,0}^{-1}-\tilde A-y\tilde A')
        \mathbf{e}'_{k-i}\pmod{y^{-t+2}}
    \end{align*}
Hence, we observe that $s_{k,i}\leq t=2k-r'$.


We let
    \begin{align*}
        g_{k,i}(T)
        &:=
        \begin{bmatrix}
            T^{k-1} & \cdots & 1
        \end{bmatrix}
        (B_{k,0}^{-1}-\tilde A)
        \mathbf{e}'_{k-i};
        \\
        h_{k,i}(T)
        &:=
        \begin{bmatrix}
            T^{k-1} & \cdots & 1
        \end{bmatrix}
        \tilde A'
        \mathbf{e}'_{k-i},
    \end{align*}
so that we have
    \begin{align}\label{eq: tilde P_k,i}
    \begin{split}
        \tilde P_{k,i}(T,x)
        &\equiv
        y^{-t}\big(1+y\Tr((A'_k)^{-1})\big)
        \big(g_{k,i}(T)-yh_{k,i}(T)\big)
        \\
        &\equiv
        y^{-t}
        g_{k,i}(T)
        +
        y^{-t+1}
        \big(
        \Tr((A'_k)^{-1})g_{k,i}(T)
        -h_{k,i}(T)
        \big)\pmod{y^{-t+2}}.
    \end{split}
    \end{align}


Consider first the case $0\leq i\leq t-1$. We see that $g_{k,i}(T)=0$ as the $(k-i)$-th column of $B_{k,0}^{-1}-\tilde A$ is zero (see \eqref{eq: inverse of B_k,x-tilde A}). As $(k-i)$-th column of $\tilde A'$ is not zero, $h_{k,i}(T)\neq 0$. So, we get $s_{k,i}=t-1$ (from \eqref{eq: tilde P_k,i}) and
    \begin{align*}
        P_{k,i,t-1}(T)
        =h_{k,i}(T).
    \end{align*}
We further characterize $h_{k,i}(T)$. Observe that, from the definition of $h_{k,i}(T)$ together with \eqref{eq: tilde A prime}, we can easily see that
    \begin{align*}
        \deg_T\big(h_{k,i}(T)-T^{i+k-t}\big)
        \leq
        k-t-1.
    \end{align*}
Also, we see that
    \begin{align*}
        [Tf_{r'}(T)h_{k,i}(T)]_{T}^{(r'-1)}
        &=
        \begin{bmatrix}
            T^{k-1} & \cdots & 1
        \end{bmatrix}
        M_{r',0}
        \begin{bmatrix}
            0_{r'-k,t} \\
            \tilde A'
        \end{bmatrix}
        \mathbf{e}'_{k-i}
        \\&=
        \begin{bmatrix}
            T^{k-1} & \cdots & 1
        \end{bmatrix}
        \begin{bmatrix}
          B_{k,0}\tilde A' \\
          D_{k,0}\tilde A'
        \end{bmatrix}
        \mathbf{e}'_{k-i}
        \\&=
        \begin{bmatrix}
            T^{k-1} & \cdots & 1
        \end{bmatrix}
        \begin{bmatrix}
            0_{k-t,k} \\
            ((A'_k)^{-1})^2V_tB_{k,0}^{-1}\\
            D_{k,0}\tilde A'
        \end{bmatrix}
        \mathbf{e}'_{k-i}
    \end{align*}
and so
    \begin{align*}
        \deg_T [Tf_{r'}(T)h_{k,j}(T)]_{T}^{(r'-1)} \leq (r'-1)-(k-t).
    \end{align*}
Noting that $k-t=r'-k$, we observe that
    \begin{align*}
        \begin{cases}
            \deg_T (h_{k,i}(T)-T^{i+r'-k})\leq r'-k-1,\\
            \deg_T [Tf_{r'}(T)h_{k,i}(T)]_{T}^{(r'-1)} \leq k-1.
        \end{cases}
    \end{align*}
According to Lemma~\ref{lemma: PQdegree}, $P_{r'-k,i+r'-k}(T,0)$ is the unique polynomial satisfying such properties.
Therefore, we have
    \begin{align*}
        P_{k,i,t-1}(T)
        =h_{k,i}(T)
        =P_{r'-k,i+r'-k}(T,0).
    \end{align*}



We now consider the case $t\leq i\leq k-1$. From the definition of $g_{k,i}(T)$ together with \eqref{eq: inverse of B_k,x-tilde A}, we see that $g_{k,i}(T)\neq 0$ and $s_{k,i}=t$. Moreover, from \eqref{eq: tilde P_k,i} we have
    \begin{align*}
        P_{k,i,t}(T) =g_{k,i}(T).
    \end{align*}
Setting $\mathbf{e}''_1,\cdots,\mathbf{e}''_{k-t}$ be the standard basis of $\Q^{k-t}$, we have
    \begin{align*}
        g_{k,i}(T)
        =
        \begin{bmatrix}
            T^{k-t-1} & \cdots & 1
        \end{bmatrix}
        B_{k-t,0}^{-1}
        \mathbf{e}''_{k-i}=P_{k-t,i-t}(T,0).
    \end{align*}
Noting that $k-t=r'-k$, we get
    \begin{align*}
        P_{k,i,t}(T)=P_{r'-k,i+r'-2k}(T,0).
    \end{align*}


We finally consider the case $k\leq i\leq r'-1$. For $k\leq i\leq r'-1$, we have
    \begin{align*}
        P_{k,i}(T,x)
        &=
        T^i
        -
        \begin{bmatrix}
            T^{k-1} & \cdots & 1
        \end{bmatrix}
        B_{k,x}^{-1}A_{k,x}
        \mathbf{e}''_{r'-i}.
    \end{align*}
Set $y=x^{-1}$ as before, then
    \begin{align*}
        A_{k,x}\mathbf{e}''_{r'-i}
        &=
        x
        \begin{bmatrix}
           I_{k-t} \\
           0_{t,k-t}
        \end{bmatrix}
        \mathbf{e}''_{r'-i}
        +
        A_{k,0}\mathbf{e}''_{r'-i}
        =
        y^{-1}\mathbf{e}'_{r'-i}
        +A_{k,0}\mathbf{e}''_{r'-i}
        \\
        &\equiv
        y^{-1}\mathbf{e}'_{r'-i}
        \pmod{y^0}.
    \end{align*}
Using $B_{k,x}^{-1}\in \M_{k\times k}(\Q[[y]])$, one get $P_{k,i-r'+k}(T,x)\in (\Q[[y]])[T]$, and
    \begin{align*}
        P_{k,i}(T,x)
        &\equiv
        -
        \begin{bmatrix}
            T^{k-1} & \cdots & 1
        \end{bmatrix}
        B_{k,x}^{-1}(y^{-1}\mathbf{e}'_{r'-i})
        \\
        &\equiv
        -y^{-1}
        \begin{bmatrix}
            T^{k-1} & \cdots & 1
        \end{bmatrix}
        B_{k,x}^{-1}\mathbf{e}'_{k-(i-r'+k)}
        \\
        &\equiv
        -y^{-1}P_{k,i-r'+k}(T,x)
        \pmod{y^0}.
    \end{align*}
So this gives
    \begin{align*}
        \frac{\tilde P_{k,i}(T,x)}{d_k(x)}
        \equiv
        -y^{-1}\frac{\tilde P_{k,i-r'+k}(T,x)}{d_k(x)}
        \pmod{y^0}.
    \end{align*}
As $d_k(x)$ is monic of degree $t$, one has
    \begin{align*}
        d_k(x)
        \equiv
        y^{-t} \pmod{y^{-t+1}}.
    \end{align*}
Also, as $t\leq i-r'+k\leq k-1$, one has
    \begin{align*}
        \tilde P_{k,i-r'+k}(T,x)
        \equiv
        y^{-t}P_{r'-k,i-k}(T,0) \pmod{y^{-t+1}}.
    \end{align*}
Combining these gives
    \begin{align*}
        \tilde P_{k,i}(T,x)
        &\equiv -y^{-1}\tilde P_{k,i-r'+k}(T,x)
        \\&\equiv
        -y^{-t-1}P_{r'-k,i-k}(T,0) \pmod{y^{-t}}.
    \end{align*}
Therefore, we get $s_{k,i}=t+1$ and
    \begin{align*}
        P_{k,i,s_{k,i}}(T)=-P_{r'-k,i-k}(T,0)
    \end{align*}
for $k\leq i\leq r'-1$.

Lastly, from $\tilde Q_{k,i}(T,x)=\big[\tilde P_{k,i}(T,x)(x+Tf_{r'}(T))\big]_{T}^{(r'-1)}$ we have
    \begin{multline*}
        \tilde Q_{k,i}(T,x)=
        x^{s_{k,i}+1} P_{k,i,s_{k,i}}(T)\\ +\sum_{s=1}^{s_{k,i}}x^s\big[P_{k,i,s}(T)Tf_{r'}(T)+P_{k,i,s-1}(T)\big]_{T}^{(r'-1)}
        +\big[P_{k,i,0}(T)Tf_{r'}(T)\big]_{T}^{(r'-1)},
    \end{multline*}
and so we see that $t_{k,i}=s_{k,i}+1$ and $Q_{k,i,t_{k,i}}(T)=P_{k,i,s_{k,i}}(T)$.
\end{proof}

\end{document}
